\documentclass[
	11pt,								% globale Schriftgröße
	parskip=half-,						% setzt Absatzabstand hoch
	paper=a4,							% Format
	english,ngerman,					% lädt Sprachpakete
	]{scrartcl}							% Dokumentenklasse
	
% //////////////////// Pakete laden ////////////////////
\usepackage{amsmath}			% MUSS vor fontspec geladen werden
\usepackage{mathtools}			% modifiziert amsmath
\usepackage{amssymb}			% mathematische symbole, für \ceckmarks
\usepackage{amsthm}				% für proof
\usepackage{mathrsfs}			% für \mathscr
\usepackage{latexsym}

\usepackage{ wasysym } % für lightning symbol

\usepackage[utf8]{inputenc}
\usepackage{microtype}			% für bessere Worttrennung
\usepackage[ngerman]{babel} 	% Spracheinstellung
\usepackage{lmodern}			% verändert verwendete Schriftart, damit sie weniger pixelig ist

\usepackage{verbatim}
\usepackage{listings}			% Für Quellcode

\usepackage{graphicx}
\usepackage{fullpage}
\usepackage{multirow}			% für multirow in tabulars
\usepackage{rotate}
\usepackage[cmyk,table]{xcolor} % um Farben zu benutzen, kann mehr als das Paket color
\usepackage[					% Verlinkungen
	colorlinks,					% farbige Schrift, statt farbiger Rahmen
	linktocpage,				% verlinkt im Abb.Verzeichnis Seitenzahl statt Bildunterschrift
	linkcolor=blue				% setzt Farbe der Links auf blau
	]{hyperref}					% nur für digitale Anwendungen, url = "http://www.example.com"
\usepackage{url}				% für Webadressen wie e-mail usw.: "\url{http://www.example.com}"

\usepackage{enumerate}			% für versch. Aufzählungezeichen wie z.B. a)
\usepackage{xspace}				% folgt ein Leerzeichen nach einem \Befehl, wird es nicht verschluckt.
\usepackage{cancel}				% für das Durchstreichen u.a. in Matheformeln mit \cancel

% zum Zeichnen (u.a. von Graphen)
\usepackage{fp}
\usepackage{tikz}
\usetikzlibrary{positioning}	% verbesserte Positionierung der Knoten
\usetikzlibrary{automata}		% für Automaten (GTI)
\usetikzlibrary{arrows}
\usetikzlibrary{shapes}
\usetikzlibrary{decorations.pathmorphing}
\usetikzlibrary{decorations.pathreplacing}
\usetikzlibrary{decorations.shapes}
\usetikzlibrary{decorations.text}

% //////////////////// Syntaxhighlighting ////////////////////
\lstloadlanguages{Python, Haskell, [LaTeX]TeX, Java}
\lstset{
   basicstyle=\footnotesize\ttfamily,	% \scriptsize the size of the fonts that are used for the code
   backgroundcolor = \color{bgcolour},	% legt Farbe der Box fest
   breakatwhitespace=false,	% sets if automatic breaks should only happen at whitespace
   breaklines=true,			% sets automatic line breaking
   captionpos=t,				% sets the caption-position to bottom, t for top
   commentstyle=\color{codeblue}\ttfamily,% comment style
   frame=single,				% adds a frame around the code
   keepspaces=true,			% keeps spaces in text, useful for keeping indentation
							% of code (possibly needs columns=flexible)
   keywordstyle=\bfseries\ttfamily\color{codepurple},% keyword style
   numbers=left,				% where to put the line-numbers;
   							% possible values are (none, left, right)
   numberstyle=\tiny\color{codegreen},	% the style that is used for the line-numbers
   numbersep=5pt,			% how far the line-numbers are from the code
   stepnumber=1,				% nummeriert nur jede i-te Zeile
   showspaces=false,			% show spaces everywhere adding particular underscores;
							% it overrides 'showstringspaces'
   showstringspaces=false,	% underline spaces within strings only
   showtabs=false,			% show tabs within strings adding particular underscores
   flexiblecolumns=false,
   tabsize=1,				% the step between two line-numbers. If 1: each line will be numbered
   stringstyle=\color{orange}\ttfamily,	% string literal style
   numberblanklines=false,				% leere Zeilen werden nicht mitnummeriert
   xleftmargin=1.2em,					% Abstand zum linken Layoutrand
   xrightmargin=0.4em,					% Abstand zum rechten Layoutrand
   aboveskip=2ex, 
}

\lstdefinestyle{py}{
   language=Python,
}
\lstdefinestyle{hs}{
   language=Haskell,
}
\lstdefinestyle{tex}{
	language=[LaTeX]TeX,
	escapeinside={\%*}{*)},     % if you want to add LaTeX within your code
	texcsstyle=*\bfseries\color{blue},% hervorhebung der tex-Schlüsselwörter
	morekeywords={*,$,\{,\},\[,\]},% if you want to add more keywords to the set
}
\lstdefinestyle{java}{
	language=Java,
	extendedchars=true,		% lets you use non-ASCII characters;
   							% for 8-bits encodings only, does not work with UTF-8
}

% //////////////////// eigene Kommandos ////////////////////
\newcommand{\Aufgabe}[2]{#1. Aufgabe \hfill \normalsize\textit{(#2 Punkte)}}
\newcommand{\Teilaufgabe}[1]{#1}
\newcommand{\Quellcode}[3]{\lstset{style=#2}% 1. path/filename, 2. type, 3. beschr.
	\lstinputlisting[caption={#3}]{#1.#2}}
\newcommand\FU{Freie Universität Berlin\xspace}% benötigt package xspace
\newcommand\gdw{g.\,d.\,w.\xspace}
\newcommand\oBdA{o.\,B.\,d.\,A.\xspace}
\newcommand{\Eu}{\texteuro}
\newcommand\N{\mathbb{N}\xspace}
\newcommand\Q{\mathbb{Q}\xspace}
\newcommand\R{\mathbb{R}\xspace}
\newcommand\Z{\mathbb{Z}\xspace}
\newcommand\ohneNull{\ensuremath{\backslash\lbrace 0\rbrace}}% \{0}
\let\dhALT\dh	% Schreibt Befehl \dh in \dhALT um
\renewcommand\dh{d.\,h.\xspace}	%renew überschreibt command \dh
\newcommand\Bolt{\;\text{\LARGE\raisebox{-0.3em}{\Lightning}\normalsize}\xspace}% Blitz
\newcommand\zz{\ensuremath{\raisebox{+0.25ex}{Z}% zu zeigen
			\kern-0.4em\raisebox{-0.25ex}{Z}%
			\;\xspace}}
\newcommand{\from}{\ensuremath{\colon}}
\newcommand{\floor}[1]{\lfloor{#1}\rfloor}
\newcommand{\ceil}[1]{\lceil{#1}\rceil}
% %\renewcommand{\L}{\ensuremath{\mathcal{L}}\xspace}
% %\renewcommand{\P}{\ensuremath{\mathcal{P}}\xspace}
% %\newcommand{\NL}{\ensuremath{\mathcal{N}\kern-0.2em\mathcal{L}}\xspace}
% %\newcommand{\NP}{\ensuremath{\mathcal{NP}}\xspace}

% //////////////////// Mathefunktionen ////////////////////
\DeclareMathOperator{\Landau}{\mathcal{O}}

% //////////////////// eigene Theoreme ////////////////////
\newtheorem{theorem}{Satz}
\newtheorem{corollary}[theorem]{Folgerung}
\newtheorem{lemma}[theorem]{Lemma}
\newtheorem{observation}[theorem]{Beobachtung}
\newtheorem{definition}[theorem]{Definition}
\newtheorem{Literatur}[theorem]{Literatur}
% konfiguriert proof
\makeatletter
\newenvironment{Proof}[1][\proofname]{\par
  \pushQED{\qed}%
  \normalfont \topsep6\p@\@plus6\p@\relax
  \trivlist
  \item[\hskip\labelsep
%         \itshape
        \bfseries
    #1\@addpunct{.}]\ignorespaces
}{%
  \popQED\endtrivlist\@endpefalse
}
\makeatother

% //////////////////// eigene Farben ////////////////////
\let\definecolor=\xdefinecolor
\definecolor{FUgreen}{RGB}{153,204,0}
\definecolor{FUblue}{RGB}{0,51,102}

\definecolor{middlegray}{rgb}{0.5,0.5,0.5}
\definecolor{lightgray}{rgb}{0.8,0.8,0.8}
\definecolor{orange}{rgb}{0.8,0.3,0.3}
\definecolor{azur}{rgb}{0,0.7,1}
\definecolor{yac}{rgb}{0.6,0.6,0.1}
\definecolor{Pink}{rgb}{1,0,0.6}

\definecolor{bgcolour}{rgb}{0.97,0.97,0.97}
\definecolor{codegreen}{rgb}{0,0.6,0}
\definecolor{codegray}{rgb}{0.35,0.35,0.35}
\definecolor{codepurple}{rgb}{0.58,0,0.82}
\definecolor{codeblue}{rgb}{0.4,0.5,1}

% //////////////////// eigene Settings ////////////////////

\textheight = 230mm		% Höhe des Satzspiegels / Layouts
\footskip = 10ex			% Abstand zw. Fußzeile und Grundlinie letzter Textzeile
\parindent 0pt			% verhindert Einrückung der 1. Zeile eines Absatzes
\setkomafont{sectioning}{\rmfamily\bfseries}% setzt Ü-Schriften in Serifen, {disposition}

\lstdefinelanguage{asm} {
  morekeywords={
    nop,movl,pushl,popl,cmpl,testl,leal,cmovl,jmp,je,jne,jz,jnz,jg,jge,jl,jle,addl,subl,imul,idiv,cdq,incl,decl,negl,andl,orl,xorl,notl,shrl,shll,sarl,sall,ret,leave,call,setg,setl,setge,setle,setne,sete,movzbl,eax,ebx,ecx,edx,edi,esi,ebp,esp,al,bl,cl,dl
  },
  sensitive=false,
  morecomment=[l]{\#},
  morestring=[b]"
}
						% bindet Header ein (WICHTIG)

% /////////////////////// BEGIN DOKUMENT /////////////////////////
\begin{document}

% /////////////////////// BEGIN TITLEPAGE /////////////////////////
\begin{titlepage}
	\subject{Jochen Schiller}					% <-- Hier Dozenten-Name eintragen
	\title{Rechnerarchitektur}	% <-- Hier Lehrveranst.Name eintragen
	\subtitle{\Large WS 16/17, Übung 1}	% <-- Hier Semester und Nummer eintragen
	\author{%
    	Tutor: Leonard König\\			% <-- Hier Tutor-Namen eintragen
        Tutorium: Donnerstag, 16-18 Uhr\\ \\	% <-- Hier Tutoriumzeit eintragen
        Studierende: Kathrin Le und Nils Koppelmann}		% <-- Hier eure Namen eintragen
	\date{\normalsize \today}					% aktuelles Datum mit \today
\end{titlepage}

\maketitle								% Erstellt das Titelblatt
\vspace*{-11cm}							% rückt Logo an den oberen Seitenrand
\makebox[\dimexpr\textwidth+1cm][r]{	%rechtsbündig und geht rechts 1cm über Layout hinaus
	\includegraphics[width=0.4\textwidth]{src/fu_logo} % fügt FU-Logo ein
}
% /////////////////////// END TITLEPAGE /////////////////////////

\vspace{8.2cm}							% Abstand
\rule{\linewidth}{0.8pt}				% horizontale Linie

% /////////////////////// Aufgabe 1 /////////////////////////
\section*{\Aufgabe{1}{Punkte}}

\subsection*{Interpreter}
Im Vergleich zum Compiler übersetzt der Interpreter den übergebenen Quellcode nicht in ein maschinenlesbares Executable sondern analysiert diesen und führt in aus.\footnote{\url{https://de.wikipedia.org/wiki/Interpreter}}

\subsection*{Compiler}
Ein Compiler übersetzt den übergebenen Quellcode zB. einer Hochsprache in ein maschinenlesbares Executable.\footnote{\url{https://de.wikipedia.org/wiki/Compiler}}

\subsection*{Assembler}
Ein Assembler (Zusammensetzer) übersetzt Assemblersprache in Maschinensprache.\footnote{\url{https://de.wikipedia.org/wiki/Assembler\_(Informatik)}}

\subsection*{Programmiersprache}
Eine Programmiersprache ist ein formales Werkzeug zur Implementierung von Algorithmen und Datenstrukturen. Das Ziel ist einer Rechenmaschiene verständlich Anweisungen zu geben.\footnote{\url{https://wiki.ubuntuusers.de/Programmiersprache/}}

\subsection*{Hochsprache}
Eine Hochsprache ist eine Programmiersprache die nicht hardwarenah ist und entsprechend kompiliert oder interpretiert werden muss. Sie unterscheidet sich sowohl in Abstraktion als auch Komplexität von hardwarenahen Programmiersprachen wie zb. Assemblersprachen.\footnote{\url{https://de.wikipedia.org/wiki/Höhere\_Programmiersprache}}

\subsection*{Assemblersprache}
Eine Assemblersprache ist eine hardwarenahe und oft hardwarespezifische imperative Programmiersprache, die durch einen Assembler kompiliert werden muss.\footnote{\url{https://de.wikipedia.org/wiki/Assemblersprache}}

\subsection*{ISA}
ISA steht für instruction set architecture und erlaubt die Nutzung von Programmen unabhängig von der Implementierung des jeweiligen Mikroprozessors.

\subsection*{CISC}
CISC steht für Complex Instruction Set Computing und beschreibt einen Prozessor mit komplexem Befehlsatz.

\subsection*{RISC}
RISC steht für Reduced Instruction Set Computing und ist im Vergleich zu CISC ein Prozessor mit einem auf die elementaren Befehle beschränkter Befehlsatz.\footnote{\url{http://www.elektronik-kompendium.de/sites/com/0412281.htm}}


% /////////////////////// END DOKUMENT /////////////////////////
\end{document}
